% *==================================================================================*
% *                     Review vs. Camera-Ready settings                             *
% *==================================================================================*
%
% REVIEW: Use the following command for submitting the paper (double-blind,
% for review):
% \documentclass{Interspeech}
%
% CAMERA-READY: Use the following command for the camera-ready version, one
% affiliation per line:
\documentclass[cameraready]{Interspeech}
% *==================================================================================*

\usepackage{multirow}
\usepackage{booktabs}
\usepackage{amssymb}

%==================================================================================
% Title
% Must exactly match the title entered into the paper submission system
\title{Dual-Space Constrained Face-Based Zero-Shot Text-to-Speech Synthesis}

%==================================================================================
% Authors
% The order of authors here must exactly match the order entered into the paper submission system
% Note that the COMPLETE list of authors MUST be entered into the paper submission system at the outset, including when submitting your manuscript for double-blind review
% The ORCID number is still optional but will become mandatory in the future years. It is strongly encouraged to get an ORCID for each cu-author.
% Middle names, including initials, must be included in the first name
\author[affiliation={1}]{Jianrong}{Wang}
\author[affiliation={1}]{Shengjie}{Zhou}
\author[affiliation={2}, correspondingauthor]{Ju}{Zhang}
\author[affiliation={1}]{Dengcheng}{Hu}
\author[affiliation={3}]{Qi}{Li}
% The maximum number of authors in the author list is 20. If the number of contributing authors is more than this, they should be listed in a footnote or the acknowledgement section.

%==================================================================================
% Affiliations

\address{
    $^1$ College of Intelligence and Computing, Tianjin University, Tianjin 300350, China \\
    $^2$ Technical College for the Deaf, Tianjin University of Technology, Tianjin 300384, China \\
    $^3$ School of Electrical and Information Engineering, Tianjin University, Tianjin 300072, China
}

%==================================================================================
% Emails
\email{wjr@tju.edu.cn, zhoushengjie@tju.edu.cn, juzhang@email.tjut.edu.cn, hdc@tju.edu.cn, qi.li@tju.edu.cn}

%==================================================================================
% Keywords
\keywords{face-based TTS, face-voice alignment, zero-shot TTS, speaker consistency}

%==================================================================================
% Content

\begin{document}

\maketitle

\begin{abstract}

A human face conveys rich cues about speaker identity, enabling face-based zero-shot text-to-speech (TTS) for unseen speakers. However, in modular face-based TTS systems, the acoustic model is typically trained on speech-derived embeddings, while face-derived representations are introduced only at inference time, often resulting in identity drift. We propose Dual-Space Constrained TTS (DSC-TTS), a modular framework that enforces identity consistency during acoustic model training in both the speaker embedding space and a shared identity space learned through face-voice alignment. By constraining representations across these complementary spaces, the proposed framework improves speaker identity stability while preserving speech quality. Experiments demonstrate higher speaker similarity and stronger identity consistency than existing face-based TTS methods.

\end{abstract}

\section{Introduction}

Generating speaker-consistent speech from a single face image is a central goal of face-based zero-shot text-to-speech (TTS), as it enables speech synthesis for previously unseen individuals without requiring enrollment audio. This capability is particularly valuable in scenarios where reference speech is unavailable, such as dubbing for silent films, virtual avatar animation, and face-based human--computer interaction. Despite recent progress, stable speaker identity from a single face image remains difficult to achieve. In the zero-shot setting, the absence of enrollment audio makes it difficult to reliably ground speaker embeddings, a problem further aggravated by heterogeneous datasets and domain shifts. As a result, high similarity between face-derived embeddings does not necessarily lead to perceptually consistent speaker identity in synthesized speech, often leading to unstable identity and timbre inconsistency. This gap highlights a fundamental challenge in face-based TTS: speaker identity must be inferred and preserved without any acoustic reference for speaker grounding.

Existing face-based TTS methods can be broadly divided into two paradigms. End-to-end approaches integrate facial representations directly into acoustic model training to improve speaker consistency~\cite{10094745,lee24_interspeech,jeon-etal-2025-progressive}. However, these methods are often limited by the scale and quality of available audio--visual datasets, and they typically fall short of the naturalness and speaker fidelity achieved by TTS models trained on large-scale speech-only corpora. Modular approaches, in contrast, decouple cross-modal representation learning from acoustic modeling, using face-derived representations to predict speaker conditioning inputs for pretrained multi-speaker TTS models. Face2Speech~\cite{goto20_interspeech} pioneered this paradigm, followed by work on improved alignment objectives~\cite{9687866}, residual-guided supervision~\cite{9746808}, structure-preserving cross-modal alignment~\cite{10448433}, and disentanglement of speaker identity from prosody~\cite{10889667}. While modular systems benefit from strong speech synthesis backbones, they remain prone to speaker inconsistency when conditioned on face-derived representations.

Several factors contribute to this limitation in modular face-based TTS systems. Speaker representations learned from heterogeneous corpora often exhibit corpus-dependent variation, which undermines identity stability when transferred across datasets. In addition, face--voice alignment is typically learned in a weak or unidirectional manner, making it sensitive to modality mismatch and domain shift. More importantly, face-derived conditioning is usually introduced only at inference time, whereas the acoustic model is trained exclusively with speech-derived speaker representations. This mismatch between training and inference representations further exacerbates identity inconsistency during face-conditioned synthesis.

To address these challenges, we propose a three-stage modular framework termed \emph{Dual-Space Constrained TTS} (DSC-TTS). The key idea is to constrain speaker identity representations in two complementary spaces: a speaker embedding space learned from speech and a shared identity space learned through face--voice alignment. Specifically, we first learn speaker embeddings that reduce corpus-dependent variation while preserving speaker-discriminative structure. We then train a symmetric, bidirectional face--voice alignment module to map facial and speech representations into a shared latent space, enabling cross-modal identity consistency. Finally, during acoustic model training, identity representations are jointly constrained in both spaces, which improves compatibility with face-derived conditioning at inference time and mitigates identity drift in synthesized speech. Audio samples are available on our project page.\footnote{\url{https://zsj23.github.io/dsc-tts}}

\begin{figure*}[t]
  \centering
  \includegraphics[width=\textwidth]{framework.png}
  \caption{Overview of the proposed DSC-TTS framework.}
  \label{fig:overview}
\end{figure*}

The main contributions of this work are as follows:
\begin{itemize}
    \item We propose a speaker embedding learning strategy that reduces corpus-dependent variation while preserving speaker-discriminative structure for modular face-based TTS.
    \item We design a symmetric, bidirectional face--voice alignment method to learn a shared identity space that enhances cross-modal consistency.
    \item We introduce DSC-TTS, a modular framework that enforces consistency in both the speaker embedding space and the face--voice shared identity space, thereby mitigating training--inference mismatch and reducing identity drift.
\end{itemize}

\section{Method}

We propose DSC-TTS, a three-stage modular framework for face-based zero-shot TTS, consisting of a speech encoder, a face--voice alignment module, and a dual-space constrained training strategy for the TTS acoustic model. As shown in Fig.~\ref{fig:overview}, the speech encoder and the face--voice alignment module are trained in the first two stages. In the final stage, the learned identity representations are used to constrain acoustic model training.

\subsection{Speech Encoder}
\label{sec:speech_encoder}

The speech encoder $E_s$ extracts fixed-dimensional, $\ell_2$-normalized speaker embeddings from log-Mel spectrograms and is trained using speaker classification as the primary objective, augmented with auxiliary objectives for corpus invariance and speaker-discriminative structure preservation.

Speaker discrimination is optimized using AAM-Softmax~\cite{8953658}. To mitigate corpus-dependent variation, we adopt domain-adversarial training via a gradient reversal layer~\cite{JMLR:v17:15-239}. However, enforcing domain invariance may weaken speaker-discriminative structure. To counteract this effect, we apply supervised contrastive learning (SupCon)~\cite{NEURIPS2020_d89a66c7}, which promotes intra-speaker compactness and inter-speaker separability.

The overall training objective is
\begin{equation}
\mathcal{L}_{\mathrm{SE}}
=
\mathcal{L}_{\mathrm{AAM}}
+
\lambda_{\mathrm{Domain}} \mathcal{L}_{\mathrm{Domain}}
+
\lambda_{\mathrm{SupCon}} \mathcal{L}_{\mathrm{SupCon}},
\end{equation}
where $\lambda_{\mathrm{Domain}}$ and $\lambda_{\mathrm{SupCon}}$ control the contribution of the auxiliary losses. After training, $E_s$ is fixed and reused in subsequent stages.

\subsection{Face--Voice Latent Alignment}

For each video, $K$ face frames are uniformly sampled and aggregated into a single face embedding $\mathbf{f}$ using attention pooling. The corresponding audio is encoded by the speech encoder $E_s$ to obtain a speaker embedding $\mathbf{v}$. Face and voice embeddings are mapped into a shared latent space through modality-specific projection networks,
\begin{align}
\mathbf{z}_f &= P_f(\mathbf{f}), &
\mathbf{z}_v &= P_v(\mathbf{v}),
\end{align}
and decoded back to their original embedding spaces through decoders $D_f$ and $D_v$, where all projection and decoding networks are implemented as two-layer residual MLPs, with layer normalization applied in the decoders.

The alignment module is trained using a combination of contrastive and consistency-based objectives.
Specifically, a CLIP-style contrastive loss~\cite{pmlr-v139-radford21a} is employed to bring paired face and speech representations closer in the shared latent space.
For a mini-batch of $N$ paired samples $\{(\mathbf{z}_f^i,\mathbf{z}_v^i)\}_{i=1}^{N}$, the InfoNCE objective is defined as
\begin{equation}
\mathcal{L}_{\mathrm{CLIP}}
=
-\frac{1}{N}
\sum_{i=1}^{N}
\log
\frac{
\exp\bigl(\cos(\mathbf{z}_f^i,\mathbf{z}_v^i)/\tau\bigr)
}{
\sum_{j=1}^{N}
\exp\bigl(\cos(\mathbf{z}_f^i,\mathbf{z}_v^j)/\tau\bigr)
},
\end{equation}
and is applied symmetrically in both face-to-voice and voice-to-face directions.

Using cosine distance $d(a,b)=1-\cos(a,b)$, the reconstruction loss supervises decoding from the shared latent space,
\begin{equation}
\mathcal{L}_{\mathrm{REC}}
=
d\Bigl(D_v(\mathbf{z}_f), \mathbf{v}\Bigr)
+
d\Bigl(D_f(\mathbf{z}_v), \mathbf{f}_{\mathrm{avg}}\Bigr),
\end{equation}
where $\mathbf{f}_{\mathrm{avg}}$ denotes the average-pooled face embedding serving as a fixed reconstruction target.

To further regularize the decoders, we introduce a latent consistency loss,
\begin{equation}
\mathcal{L}_{\mathrm{LC}}
=
d\Bigl(\mathbf{z}_f, P_v(D_v(\mathbf{z}_f))\Bigr)
+
d\Bigl(\mathbf{z}_v, P_f(D_f(\mathbf{z}_v))\Bigr),
\end{equation}
where the decoded embeddings are detached to prevent shortcut solutions.

The overall alignment objective is
\begin{equation}
\mathcal{L}_{\mathrm{ALIGN}}
=
\mathcal{L}_{\mathrm{CLIP}}
+
\lambda_{\mathrm{REC}} \mathcal{L}_{\mathrm{REC}}
+
\lambda_{\mathrm{LC}} \mathcal{L}_{\mathrm{LC}}.
\end{equation}

\begin{table*}[t]
\caption{
Comparison with face-based TTS methods on VoxCeleb2 and LRS2. For modular baselines, we re-implement face--voice alignment architectures and objectives from prior work within a unified TTS framework. All systems adopt the same acoustic model and speaker encoder architecture, with differences arising from alignment design and speaker embedding learning strategies.
}
\label{tab:comparison}
\centering
\small
\setlength{\tabcolsep}{3.2pt}
\begin{tabular*}{\textwidth}{@{\extracolsep{\fill}}ll|cccccc|cccc}
\toprule
Paradigm & Method
& \multicolumn{6}{c|}{VoxCeleb2}
& \multicolumn{4}{c}{LRS2} \\
 &
& WER$\downarrow$ & CER$\downarrow$ & SECS$\uparrow$ & SEC$\uparrow$ & SED$\downarrow$ & SMOS$\uparrow$
& WER$\downarrow$ & CER$\downarrow$ & SECS$\uparrow$ & SMOS$\uparrow$ \\
\midrule
End-to-End
& FaceTTS
& 0.113 & 0.064 & 0.603 & 0.823 & 0.854 & 2.074$\pm$1.03
& 0.119 & 0.077 & 0.552 & 2.143$\pm$1.12 \\
\midrule
\multirow{4}{*}{Modular}
& Face2Speech
& 0.098 & 0.051 & 0.613 & 0.760 & 0.842 & 2.725$\pm$1.12
& 0.077 & 0.040 & 0.569 & 2.911$\pm$1.16 \\
& SYNTHE-SEES
& 0.087 & 0.047 & 0.614 & 0.768 & 0.833 & 2.466$\pm$1.04
& 0.084 & 0.050 & 0.567 & 2.838$\pm$1.01 \\
& Face-StyleSpeech
& 0.087 & 0.044 & 0.635 & 0.744 & 0.848 & 2.658$\pm$1.07
& 0.091 & 0.050 & 0.583 & 3.066$\pm$1.00 \\
& \textbf{DSC-TTS (Ours)}
& \textbf{0.072} & \textbf{0.037} & \textbf{0.677} & \textbf{0.842} & \textbf{0.748} & \textbf{3.172$\pm$1.06}
& \textbf{0.057} & \textbf{0.032} & \textbf{0.653} & \textbf{3.416$\pm$0.92} \\
\bottomrule
\end{tabular*}
\end{table*}

\begin{table*}[t]
\caption{Ablation study on speaker embedding learning and dual-space constraints in DSC-TTS. The first row denotes the baseline trained with the standard TTS objective only, without auxiliary speaker embedding objectives or dual-space consistency constraints.}
\label{tab:ablation}
\centering
\small
\setlength{\tabcolsep}{3.2pt}
\begin{tabular*}{\textwidth}{@{\extracolsep{\fill}}cccc|ccccc|ccc}
\toprule
\multicolumn{2}{c|}{Embedding Losses}
& \multicolumn{2}{c|}{Dual-Space Losses}
& \multicolumn{5}{c|}{VoxCeleb2}
& \multicolumn{3}{c}{LRS2} \\
$\mathcal{L}_{\mathrm{Domain}}$ & $\mathcal{L}_{\mathrm{SupCon}}$
& $\mathcal{L}_{\mathrm{SCL}}$ & $\mathcal{L}_{\mathrm{SICL}}$
& WER$\downarrow$ & CER$\downarrow$ & SECS$\uparrow$ & SEC$\uparrow$ & SED$\downarrow$
& WER$\downarrow$ & CER$\downarrow$ & SECS$\uparrow$ \\
\midrule
 &  &  &
& 0.085 & 0.046 & 0.664 & 0.826 & 0.795
& 0.067 & 0.037 & 0.603 \\
$\checkmark$ &  &  &
& 0.079 & 0.041 & 0.643 & 0.838 & 0.750
& 0.073 & 0.042 & 0.624 \\
$\checkmark$ & $\checkmark$ &  &
& 0.082 & 0.044 & 0.673 & 0.841 & 0.744
& 0.070 & 0.035 & 0.641 \\
$\checkmark$ & $\checkmark$ & $\checkmark$ &
& 0.081 & 0.042 & 0.669 & 0.840 & \textbf{0.736}
& 0.064 & 0.034 & 0.637 \\
$\checkmark$ & $\checkmark$ &  & $\checkmark$
& 0.075 & 0.039 & 0.664 & \textbf{0.843} & 0.751
& 0.068 & 0.036 & 0.646 \\
$\checkmark$ & $\checkmark$ & $\checkmark$ & $\checkmark$
& \textbf{0.072} & \textbf{0.037} & \textbf{0.677} & 0.842 & 0.748
& \textbf{0.057} & \textbf{0.032} & \textbf{0.653} \\
\bottomrule
\end{tabular*}
\end{table*}

\subsection{Dual-Space Constrained TTS}

In the final stage, we explicitly constrain the acoustic model to preserve speaker identity in two complementary spaces: the speech-derived speaker embedding space and the face--voice shared identity space, thereby reducing the representation mismatch between training and face-conditioned inference. The model is pretrained on a single-speaker dataset and then fine-tuned on multi-speaker data with the speech encoder $E_s$ and the voice-side projection network $P_v$ fixed to define the shared identity space.

Given ground-truth speech $x$ and synthesized speech $\hat{x}$, speaker consistency in the embedding space is enforced by
\begin{equation}
\mathcal{L}_{\mathrm{SCL}} = -\cos\bigl(E_s(\hat{x}), E_s(x)\bigr),
\end{equation}
while consistency in the shared identity space is enforced by
\begin{equation}
\mathcal{L}_{\mathrm{SICL}} = -\cos\bigl(P_v(E_s(\hat{x})), P_v(E_s(x))\bigr).
\end{equation}

The overall training objective is
\begin{equation}
\mathcal{L}_{\mathrm{TTS}}^{\mathrm{total}}
=
\mathcal{L}_{\mathrm{TTS}}
+
\lambda_{\mathrm{SCL}} \mathcal{L}_{\mathrm{SCL}}
+
\lambda_{\mathrm{SICL}} \mathcal{L}_{\mathrm{SICL}} .
\end{equation}

\section{Experiments}

\subsection{Datasets}

VoxCeleb2~\cite{chung18b_interspeech} is a large-scale audio--visual dataset collected from unconstrained online videos and is used to train both the face--voice alignment module and the speech encoder. LibriTTS~\cite{zen19_interspeech} is a multi-speaker English speech corpus with aligned transcripts; its \texttt{train-clean-100} and \texttt{train-clean-360} subsets are used for multi-speaker TTS training and for training the speech encoder. LJSpeech~\cite{ljspeech17} is used for single-speaker TTS pretraining. All speech signals are downsampled to 16~kHz.

For evaluation, 24 unseen speakers (12 male and 12 female), disjoint from the training set, are randomly selected from VoxCeleb2. For each speaker, face images are paired with text prompts sampled from the LibriTTS \texttt{test-clean} set. LRS2~\cite{8099850} is a tri-modal face-speech-text dataset collected under different recording conditions and used for sample-level evaluation under corpus mismatch.

\subsection{Implementation Details}

We use an ECAPA-TDNN~\cite{desplanques20_interspeech} backbone for speaker embedding extraction and a pretrained AntelopeV2\footnote{\url{https://github.com/deepinsight/insightface}} for face embedding extraction.
For each video, $K=5$ face frames are uniformly sampled and aggregated with attention pooling.
Face and speaker embeddings are projected into a shared latent space of 256 dimensions.

The face--voice alignment network is trained with AdamW using a learning rate of $10^{-4}$, weight decay of $10^{-6}$, and a batch size of 256.
For TTS, we use YourTTS~\cite{pmlr-v162-casanova22a}, which is pretrained on LJSpeech for 120K steps and fine-tuned on LibriTTS for 165K steps.
Loss weights are set to $\lambda_{\mathrm{Domain}}=\lambda_{\mathrm{SupCon}}=1.0$, $\lambda_{\mathrm{REC}}=0.5$, $\lambda_{\mathrm{LC}}=0.1$, $\lambda_{\mathrm{SCL}}=9.0$, and $\lambda_{\mathrm{SICL}}=3.0$.
All experiments are conducted on a single NVIDIA RTX 4090 GPU with 24~GB memory.

\subsection{Evaluation Metrics}

Speech intelligibility is evaluated using word error rate (WER) and character error rate (CER), computed from Whisper transcriptions~\cite{pmlr-v202-radford23a}. Speaker-related properties are evaluated using Resemblyzer\footnote{\url{https://github.com/resemble-ai/Resemblyzer}}, which is used to compute speaker embedding cosine similarity between synthesized speech and the corresponding reference speech (SECS), speaker embedding consistency (SEC), and speaker embedding diversity (SED). SEC measures the average intra-speaker embedding consistency across synthesized utterances, while SED reflects inter-speaker embedding separability across different speakers.

Subjective evaluation is conducted using Similarity Mean Opinion Score (SMOS) tests on a five-point scale, which evaluate perceived speaker similarity between synthesized speech and the target speaker. Each sample is rated by 10 listeners.

\begin{figure}[t]
  \centering
  \includegraphics[width=\columnwidth]{vis_attention.png}
  \caption{
    Identity-aware attention pooling over multiple face frames.
    The model selectively emphasizes identity-informative frames while suppressing visually unreliable ones.
    \textbf{Top:} frames with degraded identity cues (e.g., extreme pose or motion blur) receive lower attention weights.
    \textbf{Bottom:} frames with non-identity variations (e.g., facial expression) are weighted more uniformly.
  }
  \label{fig:attention_pooling}
\end{figure}

\subsection{Comparison with Face-Based TTS Methods}

We compare DSC-TTS with one end-to-end baseline, FaceTTS~\cite{10094745}, and three modular methods: Face2Speech~\cite{goto20_interspeech}, SYNTHE-SEES~\cite{9746808}, and Face-StyleSpeech~\cite{10889667}. For the modular baselines, face--voice alignment components are implemented following the original designs, while the remaining components are evaluated under a common setting.

Table~\ref{tab:comparison} shows that modular methods generally produce lower WER and CER than the end-to-end baseline FaceTTS on both VoxCeleb2 and LRS2. On VoxCeleb2, FaceTTS reports a WER of 0.113, whereas all modular systems achieve WERs below 0.10. Despite better intelligibility, modular baselines exhibit weaker speaker identity consistency, as reflected by their lower SEC scores.

In contrast, DSC-TTS improves speaker-related metrics without sacrificing intelligibility. On VoxCeleb2, DSC-TTS achieves a SECS of 0.677 and an SEC of 0.842, outperforming other modular approaches while maintaining comparable speaker diversity. Consistent improvements are also observed on LRS2, suggesting that the proposed dual-space constraints generalize well under corpus mismatch. These objective gains are further supported by the SMOS results, which indicate higher perceived speaker similarity.

\subsection{Ablation Study}

Table~\ref{tab:ablation} examines the effects of speaker embedding learning and dual-space consistency constraints in DSC-TTS. Incorporating domain-adversarial and supervised contrastive objectives improves speaker-related metrics, especially SECS. The gains are more pronounced on LRS2, suggesting improved generalization of the learned speaker embeddings under corpus mismatch.

Introducing dual-space consistency constraints during multi-speaker fine-tuning further improves identity consistency across synthesized utterances. In particular, SICL improves SEC by enforcing consistency in the shared identity space, and combining SICL with SCL yields additional gains in the cross-corpus setting. Meanwhile, SED remains largely unchanged, indicating that the constraints enhance identity stability without degrading inter-speaker separability. These results show that speaker embedding learning and dual-space constrained training provide complementary improvements.

\begin{figure}[t]
  \centering
  \includegraphics[width=\columnwidth]{face_voice_discrepancy.png}
  \caption{
    Distribution of face--voice embedding discrepancy on LRS2,
    measured as $1-\cos(\tilde{\mathbf{v}}_f, \mathbf{v})$,
    where $\tilde{\mathbf{v}}_f$ denotes the face-derived speaker representation
    used for TTS conditioning.
    Left-shifted curves indicate reduced face--voice identity mismatch.
  }
  \label{fig:face_voice_discrepancy}
\end{figure}

\begin{figure}[t]
  \centering
  \includegraphics[width=\columnwidth]{vis_tsne.png}
  \caption{
    Qualitative visualization of speaker embeddings extracted from synthesized speech using t-SNE.
    More compact intra-speaker clusters and clearer inter-speaker separation indicate a more structured organization of speaker identity.
    \textbf{Left:} FaceTTS. \textbf{Right:} Ours.
  }
  \label{fig:tsne_comparison}
\end{figure}

\subsection{Visualization}

Fig.~\ref{fig:attention_pooling} illustrates the behavior of attention pooling under varying visual conditions. Frames with degraded identity cues, such as extreme pose or motion blur, receive lower attention weights, whereas frames that differ mainly in non-identity factors, such as facial expression, are weighted more uniformly. This behavior indicates that the model selectively emphasizes identity-informative frames while suppressing visually unreliable frames, which is desirable for reliable face-based speaker representation.

Fig.~\ref{fig:face_voice_discrepancy} compares the distribution of face--voice embedding discrepancy across different methods under corpus mismatch. The proposed method yields a higher fraction of low-discrepancy samples, and incorporating domain-adversarial and supervised contrastive objectives further shifts the distribution toward lower values. These results suggest that improved cross-modal alignment and reduced corpus-dependent variation jointly mitigate face--voice identity mismatch.

Fig.~\ref{fig:tsne_comparison} provides a qualitative visualization of speaker embeddings extracted from synthesized speech. Compared with FaceTTS, the proposed method produces more compact intra-speaker clusters and clearer separation between different speakers. This qualitative pattern is consistent with the quantitative trends observed in SEC and SED, indicating improved identity stability and inter-speaker separability in the synthesized speech.

\section{Conclusion}

We propose DSC-TTS, a modular face-based zero-shot text-to-speech framework that mitigates speaker identity drift from a single face image. By integrating corpus-invariant speaker embedding learning, symmetric face--voice alignment, and dual-space constrained training, DSC-TTS alleviates the training--inference mismatch in speaker identity representations. Experiments on VoxCeleb2 and LRS2 demonstrate that the proposed method improves speaker similarity and identity consistency while preserving intelligibility compared with existing face-based TTS approaches. Ablation and qualitative analyses further validate the complementary contributions of each component to the overall performance.

\section{Acknowledgments}
This work was financially supported by Key R\&D Program of the Nanning Science Research and Technology Development Plan (20253057) and Tianjin Science and Technology Program, Special Project for High-Quality Development of Manufacturing Industry (25ZGZNGX00040).

\section{Generative AI Use Disclosure}
Generative AI tools were used solely for language editing and manuscript polishing to improve grammar, clarity, and readability. No generative AI tools were used in the development of the research ideas, technical content, methodology, experimental design, results, or conclusions presented in this paper.

\bibliographystyle{IEEEtran}
\bibliography{mybib}

\end{document}

%%% Local Variables:
%%% mode: LaTeX
%%% TeX-master: t
%%% End:
